
\documentclass{article}
\usepackage{colortbl}
\usepackage{makecell}
\usepackage{multirow}
\usepackage{supertabular}

\begin{document}

\newcounter{utterance}

\centering \large Interaction Transcript for game `dond', experiment `semi\_en', episode 6 with Qwen3{-}32B{-}t0.0.
\vspace{24pt}

{ \footnotesize  \setcounter{utterance}{1}
\setlength{\tabcolsep}{0pt}
\begin{supertabular}{c@{$\;$}|p{.15\linewidth}@{}p{.15\linewidth}p{.15\linewidth}p{.15\linewidth}p{.15\linewidth}p{.15\linewidth}}
    \# & $\;$A & \multicolumn{4}{c}{Game Master} & $\;\:$B\\
    \hline

    \theutterance \stepcounter{utterance}  
    & & \multicolumn{4}{p{0.6\linewidth}}{
        \cellcolor[rgb]{0.9,0.9,0.9}{
            \makecell[{{p{\linewidth}}}]{
                \texttt{\tiny{[P1$\langle$GM]}}
                \texttt{You are playing a negotiation game in which you have to agree on how to divide a set of items among you and another player.} \\
\\ 
\texttt{Rules:} \\
\texttt{(a) You and the other player are given a set of items. Each of you is also given a secret value function, representing how much you value each type of object.} \\
\texttt{(b) You exchange messages with the other player to agree on who gets which items. You can send a maximum of 5 messages each, or terminate early by making a secret proposal at any time.} \\
\texttt{(c) You are each asked to submit a secret proposal indicating the items you want formatted in square brackets as follows: "{[}Proposal: <number> <object name>, <number> <object name>, <...>{]}"} \\
\texttt{(d) If your proposals are complementary, i.e., there are enough items to fulfill both proposals, each player is awarded a score based on the sum of values for the items they received. Otherwise, both of you get zero points.} \\
\texttt{(e) Your goal is to maximize the score you receive.} \\
\\ 
\texttt{Let us start.} \\
\\ 
\texttt{The set of available items is:} \\
\\ 
\texttt{3 apples, 1 soap, 1 egg} \\
\\ 
\texttt{Your secret value function is:} \\
\\ 
\texttt{apple: 1, soap: 3, egg: 4} \\
\\ 
\texttt{IMPORTANT: Your messages, unless it is the secret proposal, are directly transmitted to the other player, so do not include any response to the rules or text announcing your message. To make a secret proposal, use the indicated format. Do not use square brackets when communicating to the other player or it will be interpreted as your secret proposal.} \\
            }
        }
    }
    & \\ \\

    \theutterance \stepcounter{utterance}  
    & \multicolumn{4}{p{0.6\linewidth}}{
        \cellcolor[rgb]{0.9,0.9,0.9}{
            \makecell[{{p{\linewidth}}}]{
                \texttt{\tiny{[P1$\rangle$GM]}}
                \texttt{<think>} \\
\texttt{Okay, let's see. I need to negotiate the division of 3 apples, 1 soap, and 1 egg with another player. My values are apple:1, soap:3, egg:4. The other player's values are secret, so I don't know what they value. The goal is to maximize my own score, but we have to make sure our proposals are complementary. If they aren't, we both get zero. So I need to figure out a way to divide the items such that both of our proposals can be satisfied.} \\
\\ 
\texttt{First, I should consider what the other player might want. Since I don't know their values, maybe I can make an educated guess. The items are apples, soap, and egg. The other player might value the egg highly, or maybe the soap. Alternatively, they might prefer multiple apples. But since I don't know, I need to think of a proposal that allows flexibility.} \\
\\ 
\texttt{My own highest value is the egg (4), then soap (3), then apple (1). So ideally, I want the egg. But if I take the egg, I need to make sure the other player can get their desired items. Let's think about possible scenarios.} \\
\\ 
\texttt{If I propose taking the egg, then the remaining items are 3 apples and 1 soap. The other player might want the soap and some apples. If they take the soap and all 3 apples, that's possible. But if their value for soap is high, maybe they would prefer that. Alternatively, if they value apples more, maybe they want all the apples. But since I don't know, I need to make a proposal that's likely to be compatible.} \\
\\ 
\texttt{Alternatively, if I take the soap and some apples, maybe that leaves the egg for the other player. But if the other player values the egg highly, they might want it. But if I take the soap, I get 3 points, which is less than the egg's 4. So maybe I should prioritize the egg.} \\
\\ 
\texttt{Another angle: the total number of items is 5 (3 apples, 1 soap, 1 egg). If I take the egg, that's 1 item. The other player can take the remaining 4 (3 apples and 1 soap). If their proposal is to take all those, then it's complementary. But if their proposal requires more than that, like maybe they want the egg too, then it's conflicting.} \\
            }
        }
    }
    & & \\ \\

    \theutterance \stepcounter{utterance}  
    & & & \multicolumn{2}{p{0.3\linewidth}}{
        \cellcolor[rgb]{0.9,0.9,0.9}{
            \makecell[{{p{\linewidth}}}]{
                \texttt{\tiny{[GM$|$GM]}}
                \texttt{continue} \\
            }
        }
    }
    & & \\ \\

    \theutterance \stepcounter{utterance}  
    & & & \multicolumn{2}{p{0.3\linewidth}}{
        \cellcolor[rgb]{0.9,0.9,0.9}{
            \makecell[{{p{\linewidth}}}]{
                \texttt{\tiny{[GM$|$GM]}}
                \texttt{<think>} \\
\texttt{Okay, let's see. I need to negotiate the division of 3 apples, 1 soap, and 1 egg with another player. My values are apple:1, soap:3, egg:4. The other player's values are secret, so I don't know what they value. The goal is to maximize my own score, but we have to make sure our proposals are complementary. If they aren't, we both get zero. So I need to figure out a way to divide the items such that both of our proposals can be satisfied.} \\
\\ 
\texttt{First, I should consider what the other player might want. Since I don't know their values, maybe I can make an educated guess. The items are apples, soap, and egg. The other player might value the egg highly, or maybe the soap. Alternatively, they might prefer multiple apples. But since I don't know, I need to think of a proposal that allows flexibility.} \\
\\ 
\texttt{My own highest value is the egg (4), then soap (3), then apple (1). So ideally, I want the egg. But if I take the egg, I need to make sure the other player can get their desired items. Let's think about possible scenarios.} \\
\\ 
\texttt{If I propose taking the egg, then the remaining items are 3 apples and 1 soap. The other player might want the soap and some apples. If they take the soap and all 3 apples, that's possible. But if their value for soap is high, maybe they would prefer that. Alternatively, if they value apples more, maybe they want all the apples. But since I don't know, I need to make a proposal that's likely to be compatible.} \\
\\ 
\texttt{Alternatively, if I take the soap and some apples, maybe that leaves the egg for the other player. But if the other player values the egg highly, they might want it. But if I take the soap, I get 3 points, which is less than the egg's 4. So maybe I should prioritize the egg.} \\
\\ 
\texttt{Another angle: the total number of items is 5 (3 apples, 1 soap, 1 egg). If I take the egg, that's 1 item. The other player can take the remaining 4 (3 apples and 1 soap). If their proposal is to take all those, then it's complementary. But if their proposal requires more than that, like maybe they want the egg too, then it's conflicting.} \\
            }
        }
    }
    & & \\ \\

    \theutterance \stepcounter{utterance}  
    & & \multicolumn{4}{p{0.6\linewidth}}{
        \cellcolor[rgb]{0.9,0.9,0.9}{
            \makecell[{{p{\linewidth}}}]{
                \texttt{\tiny{[GM$\rangle$P2]}}
                \texttt{You are playing a negotiation game in which you have to agree on how to divide a set of items among you and another player.} \\
\\ 
\texttt{Rules:} \\
\texttt{(a) You and the other player are given a set of items. Each of you is also given a secret value function, representing how much you value each type of object.} \\
\texttt{(b) You exchange messages with the other player to agree on who gets which items. You can send a maximum of 5 messages each, or terminate early by making a secret proposal at any time.} \\
\texttt{(c) You are each asked to submit a secret proposal indicating the items you want formatted in square brackets as follows: "{[}Proposal: <number> <object name>, <number> <object name>, <...>{]}"} \\
\texttt{(d) If your proposals are complementary, i.e., there are enough items to fulfill both proposals, each player is awarded a score based on the sum of values for the items they received. Otherwise, both of you get zero points.} \\
\texttt{(e) Your goal is to maximize the score you receive.} \\
\\ 
\texttt{Let us start.} \\
\\ 
\texttt{The set of available items is:} \\
\\ 
\texttt{3 apples, 1 soap, 1 egg} \\
\\ 
\texttt{Your secret value function is:} \\
\\ 
\texttt{apple: 0, soap: 7, egg: 3} \\
\\ 
\texttt{IMPORTANT: Your messages, unless it is the secret proposal, are directly transmitted to the other player, so do not include any response to the rules or text announcing your message. To make a secret proposal, use the indicated format. Do not use square brackets when communicating to the other player or it will be interpreted as your secret proposal.} \\
\\ 
\\ 
\texttt{<think>} \\
\texttt{Okay, let's see. I need to negotiate the division of 3 apples, 1 soap, and 1 egg with another player. My values are apple:1, soap:3, egg:4. The other player's values are secret, so I don't know what they value. The goal is to maximize my own score, but we have to make sure our proposals are complementary. If they aren't, we both get zero. So I need to figure out a way to divide the items such that both of our proposals can be satisfied.} \\
\\ 
\texttt{First, I should consider what the other player might want. Since I don't know their values, maybe I can make an educated guess. The items are apples, soap, and egg. The other player might value the egg highly, or maybe the soap. Alternatively, they might prefer multiple apples. But since I don't know, I need to think of a proposal that allows flexibility.} \\
\\ 
\texttt{My own highest value is the egg (4), then soap (3), then apple (1). So ideally, I want the egg. But if I take the egg, I need to make sure the other player can get their desired items. Let's think about possible scenarios.} \\
\\ 
\texttt{If I propose taking the egg, then the remaining items are 3 apples and 1 soap. The other player might want the soap and some apples. If they take the soap and all 3 apples, that's possible. But if their value for soap is high, maybe they would prefer that. Alternatively, if they value apples more, maybe they want all the apples. But since I don't know, I need to make a proposal that's likely to be compatible.} \\
\\ 
\texttt{Alternatively, if I take the soap and some apples, maybe that leaves the egg for the other player. But if the other player values the egg highly, they might want it. But if I take the soap, I get 3 points, which is less than the egg's 4. So maybe I should prioritize the egg.} \\
\\ 
\texttt{Another angle: the total number of items is 5 (3 apples, 1 soap, 1 egg). If I take the egg, that's 1 item. The other player can take the remaining 4 (3 apples and 1 soap). If their proposal is to take all those, then it's complementary. But if their proposal requires more than that, like maybe they want the egg too, then it's conflicting.} \\
            }
        }
    }
    & \\ \\

    \theutterance \stepcounter{utterance}  
    & & & \multicolumn{4}{p{0.6\linewidth}}{
        \cellcolor[rgb]{0.9,0.9,0.9}{
            \makecell[{{p{\linewidth}}}]{
                \texttt{\tiny{[GM$\langle$P2]}}
                \texttt{<think>} \\
\texttt{Okay, let's try to figure out the best way to approach this. My values are apple:1, soap:3, egg:4. The other player's values are unknown. The key is to make a proposal that maximizes my points while ensuring that the other player's proposal doesn't conflict with mine.} \\
\\ 
\texttt{First, I need to consider what the other player might want. Since I don't know their values, I have to make assumptions. The items available are 3 apples, 1 soap, 1 egg. The other player could value any of these. If I take the egg (which I value highest), that leaves 3 apples and 1 soap. The other player might want the soap and some apples. If they take the soap and all 3 apples, that's possible. But if they also want the egg, there's a conflict. However, since I can't know their preferences, I need to structure my proposal in a way that allows the other player to take what's left without overlapping.} \\
\\ 
\texttt{Alternatively, maybe I should take the soap and some apples. If I take the soap (3 points) and, say, 1 apple (1 point), that's 4 points total. Then the other player could take the remaining 2 apples and the egg. But if the other player values the egg highly, they might want it, which would conflict if I take the soap. Wait, but if I take the soap and 1 apple, and they take the egg and 2 apples, that's possible. But my total would be 3+1=4, which is better than just taking the egg (4) alone. Hmm, same points. But maybe taking the egg is better because it's a single item, leaving more flexibility for the other player.} \\
\\ 
\texttt{Another angle: if I take the egg, the other player can take the soap and all 3 apples. That would give me 4 points, and the other player would get their value for soap and apples. If their value for soap is high, that's good for them. If they value apples more, maybe they prefer more apples. Either way, as long as their proposal doesn't require the egg, it's okay.} \\
\\ 
\texttt{But what if the other player also wants the egg? Then our proposals would conflict. But since I don't know their values, I have to assume they might not. Maybe the best strategy is to take the highest value item for} \\
            }
        }
    }
     \\ \\

    \theutterance \stepcounter{utterance}  
    & & & \multicolumn{2}{p{0.3\linewidth}}{
        \cellcolor[rgb]{0.9,0.9,0.9}{
            \makecell[{{p{\linewidth}}}]{
                \texttt{\tiny{[GM$|$GM]}}
                \texttt{continue} \\
            }
        }
    }
    & & \\ \\

    \theutterance \stepcounter{utterance}  
    & & & \multicolumn{2}{p{0.3\linewidth}}{
        \cellcolor[rgb]{0.9,0.9,0.9}{
            \makecell[{{p{\linewidth}}}]{
                \texttt{\tiny{[GM$|$GM]}}
                \texttt{<think>} \\
\texttt{Okay, let's try to figure out the best way to approach this. My values are apple:1, soap:3, egg:4. The other player's values are unknown. The key is to make a proposal that maximizes my points while ensuring that the other player's proposal doesn't conflict with mine.} \\
\\ 
\texttt{First, I need to consider what the other player might want. Since I don't know their values, I have to make assumptions. The items available are 3 apples, 1 soap, 1 egg. The other player could value any of these. If I take the egg (which I value highest), that leaves 3 apples and 1 soap. The other player might want the soap and some apples. If they take the soap and all 3 apples, that's possible. But if they also want the egg, there's a conflict. However, since I can't know their preferences, I need to structure my proposal in a way that allows the other player to take what's left without overlapping.} \\
\\ 
\texttt{Alternatively, maybe I should take the soap and some apples. If I take the soap (3 points) and, say, 1 apple (1 point), that's 4 points total. Then the other player could take the remaining 2 apples and the egg. But if the other player values the egg highly, they might want it, which would conflict if I take the soap. Wait, but if I take the soap and 1 apple, and they take the egg and 2 apples, that's possible. But my total would be 3+1=4, which is better than just taking the egg (4) alone. Hmm, same points. But maybe taking the egg is better because it's a single item, leaving more flexibility for the other player.} \\
\\ 
\texttt{Another angle: if I take the egg, the other player can take the soap and all 3 apples. That would give me 4 points, and the other player would get their value for soap and apples. If their value for soap is high, that's good for them. If they value apples more, maybe they prefer more apples. Either way, as long as their proposal doesn't require the egg, it's okay.} \\
\\ 
\texttt{But what if the other player also wants the egg? Then our proposals would conflict. But since I don't know their values, I have to assume they might not. Maybe the best strategy is to take the highest value item for} \\
            }
        }
    }
    & & \\ \\

    \theutterance \stepcounter{utterance}  
    & & \multicolumn{4}{p{0.6\linewidth}}{
        \cellcolor[rgb]{0.9,0.9,0.9}{
            \makecell[{{p{\linewidth}}}]{
                \texttt{\tiny{[P1$\langle$GM]}}
                \texttt{<think>} \\
\texttt{Okay, let's try to figure out the best way to approach this. My values are apple:1, soap:3, egg:4. The other player's values are unknown. The key is to make a proposal that maximizes my points while ensuring that the other player's proposal doesn't conflict with mine.} \\
\\ 
\texttt{First, I need to consider what the other player might want. Since I don't know their values, I have to make assumptions. The items available are 3 apples, 1 soap, 1 egg. The other player could value any of these. If I take the egg (which I value highest), that leaves 3 apples and 1 soap. The other player might want the soap and some apples. If they take the soap and all 3 apples, that's possible. But if they also want the egg, there's a conflict. However, since I can't know their preferences, I need to structure my proposal in a way that allows the other player to take what's left without overlapping.} \\
\\ 
\texttt{Alternatively, maybe I should take the soap and some apples. If I take the soap (3 points) and, say, 1 apple (1 point), that's 4 points total. Then the other player could take the remaining 2 apples and the egg. But if the other player values the egg highly, they might want it, which would conflict if I take the soap. Wait, but if I take the soap and 1 apple, and they take the egg and 2 apples, that's possible. But my total would be 3+1=4, which is better than just taking the egg (4) alone. Hmm, same points. But maybe taking the egg is better because it's a single item, leaving more flexibility for the other player.} \\
\\ 
\texttt{Another angle: if I take the egg, the other player can take the soap and all 3 apples. That would give me 4 points, and the other player would get their value for soap and apples. If their value for soap is high, that's good for them. If they value apples more, maybe they prefer more apples. Either way, as long as their proposal doesn't require the egg, it's okay.} \\
\\ 
\texttt{But what if the other player also wants the egg? Then our proposals would conflict. But since I don't know their values, I have to assume they might not. Maybe the best strategy is to take the highest value item for} \\
            }
        }
    }
    & \\ \\

    \theutterance \stepcounter{utterance}  
    & \multicolumn{4}{p{0.6\linewidth}}{
        \cellcolor[rgb]{0.9,0.9,0.9}{
            \makecell[{{p{\linewidth}}}]{
                \texttt{\tiny{[P1$\rangle$GM]}}
                \texttt{<think>} \\
\texttt{Okay, let's try to figure out the best way to approach this. My values are apple:1, soap:3, egg:4. The other player's values are unknown. The key is to make a proposal that maximizes my points while ensuring that the other player's proposal doesn't conflict with mine.} \\
\\ 
\texttt{First, I need to consider what the other player might want. Since I don't know their values, I have to make assumptions. The items available are 3 apples, 1 soap, 1 egg. The other player could value any of these. If I take the egg (which I value highest), that leaves 3 apples and 1 soap. The other player might want the soap and some apples. If they take the soap and all 3 apples, that's possible. But if they also want the egg, there's a conflict. However, since I can't know their preferences, I need to structure my proposal in a way that allows the other player to take what's left without overlapping.} \\
\\ 
\texttt{Alternatively, maybe I should take the soap and some apples. If I take the soap (3 points) and, say, 1 apple (1 point), that's 4 points total. Then the other player could take the remaining 2 apples and the egg. But if the other player values the egg highly, they might want it, which would conflict if I take the soap. Wait, but if I take the soap and 1 apple, and they take the egg and 2 apples, that's possible. But my total would be 3+1=4, which is better than just taking the egg (4) alone. Hmm, same points. But maybe taking the egg is better because it's a single item, leaving more flexibility for the other player.} \\
\\ 
\texttt{Another angle: if I take the egg, the other player can take the soap and all 3 apples. That would give me 4 points, and the other player would get their value for soap and apples. If their value for soap is high, that's good for them. If they value apples more, maybe they prefer more apples. Either way, as long as their proposal doesn't require the egg, it's okay.} \\
\\ 
\texttt{But what if the other player also wants the egg? Then our proposals would conflict. But since I don't know their values, I have to assume they might not. Maybe the best strategy is to take the highest value item for} \\
            }
        }
    }
    & & \\ \\

    \theutterance \stepcounter{utterance}  
    & & & \multicolumn{2}{p{0.3\linewidth}}{
        \cellcolor[rgb]{0.9,0.9,0.9}{
            \makecell[{{p{\linewidth}}}]{
                \texttt{\tiny{[GM$|$GM]}}
                \texttt{continue} \\
            }
        }
    }
    & & \\ \\

    \theutterance \stepcounter{utterance}  
    & & & \multicolumn{2}{p{0.3\linewidth}}{
        \cellcolor[rgb]{0.9,0.9,0.9}{
            \makecell[{{p{\linewidth}}}]{
                \texttt{\tiny{[GM$|$GM]}}
                \texttt{<think>} \\
\texttt{Okay, let's try to figure out the best way to approach this. My values are apple:1, soap:3, egg:4. The other player's values are unknown. The key is to make a proposal that maximizes my points while ensuring that the other player's proposal doesn't conflict with mine.} \\
\\ 
\texttt{First, I need to consider what the other player might want. Since I don't know their values, I have to make assumptions. The items available are 3 apples, 1 soap, 1 egg. The other player could value any of these. If I take the egg (which I value highest), that leaves 3 apples and 1 soap. The other player might want the soap and some apples. If they take the soap and all 3 apples, that's possible. But if they also want the egg, there's a conflict. However, since I can't know their preferences, I need to structure my proposal in a way that allows the other player to take what's left without overlapping.} \\
\\ 
\texttt{Alternatively, maybe I should take the soap and some apples. If I take the soap (3 points) and, say, 1 apple (1 point), that's 4 points total. Then the other player could take the remaining 2 apples and the egg. But if the other player values the egg highly, they might want it, which would conflict if I take the soap. Wait, but if I take the soap and 1 apple, and they take the egg and 2 apples, that's possible. But my total would be 3+1=4, which is better than just taking the egg (4) alone. Hmm, same points. But maybe taking the egg is better because it's a single item, leaving more flexibility for the other player.} \\
\\ 
\texttt{Another angle: if I take the egg, the other player can take the soap and all 3 apples. That would give me 4 points, and the other player would get their value for soap and apples. If their value for soap is high, that's good for them. If they value apples more, maybe they prefer more apples. Either way, as long as their proposal doesn't require the egg, it's okay.} \\
\\ 
\texttt{But what if the other player also wants the egg? Then our proposals would conflict. But since I don't know their values, I have to assume they might not. Maybe the best strategy is to take the highest value item for} \\
            }
        }
    }
    & & \\ \\

    \theutterance \stepcounter{utterance}  
    & & \multicolumn{4}{p{0.6\linewidth}}{
        \cellcolor[rgb]{0.9,0.9,0.9}{
            \makecell[{{p{\linewidth}}}]{
                \texttt{\tiny{[GM$\rangle$P2]}}
                \texttt{<think>} \\
\texttt{Okay, let's try to figure out the best way to approach this. My values are apple:1, soap:3, egg:4. The other player's values are unknown. The key is to make a proposal that maximizes my points while ensuring that the other player's proposal doesn't conflict with mine.} \\
\\ 
\texttt{First, I need to consider what the other player might want. Since I don't know their values, I have to make assumptions. The items available are 3 apples, 1 soap, 1 egg. The other player could value any of these. If I take the egg (which I value highest), that leaves 3 apples and 1 soap. The other player might want the soap and some apples. If they take the soap and all 3 apples, that's possible. But if they also want the egg, there's a conflict. However, since I can't know their preferences, I need to structure my proposal in a way that allows the other player to take what's left without overlapping.} \\
\\ 
\texttt{Alternatively, maybe I should take the soap and some apples. If I take the soap (3 points) and, say, 1 apple (1 point), that's 4 points total. Then the other player could take the remaining 2 apples and the egg. But if the other player values the egg highly, they might want it, which would conflict if I take the soap. Wait, but if I take the soap and 1 apple, and they take the egg and 2 apples, that's possible. But my total would be 3+1=4, which is better than just taking the egg (4) alone. Hmm, same points. But maybe taking the egg is better because it's a single item, leaving more flexibility for the other player.} \\
\\ 
\texttt{Another angle: if I take the egg, the other player can take the soap and all 3 apples. That would give me 4 points, and the other player would get their value for soap and apples. If their value for soap is high, that's good for them. If they value apples more, maybe they prefer more apples. Either way, as long as their proposal doesn't require the egg, it's okay.} \\
\\ 
\texttt{But what if the other player also wants the egg? Then our proposals would conflict. But since I don't know their values, I have to assume they might not. Maybe the best strategy is to take the highest value item for} \\
            }
        }
    }
    & \\ \\

    \theutterance \stepcounter{utterance}  
    & & & \multicolumn{4}{p{0.6\linewidth}}{
        \cellcolor[rgb]{0.9,0.9,0.9}{
            \makecell[{{p{\linewidth}}}]{
                \texttt{\tiny{[GM$\langle$P2]}}
                \texttt{<think>} \\
\texttt{Okay, let's try to figure out the best way to approach this. My values are apple:1, soap:3, egg:4. The other player's values are unknown. The key is to make a proposal that maximizes my points while ensuring that the other player's proposal doesn't conflict with mine.} \\
\\ 
\texttt{First, I need to consider what the other player might want. Since I don't know their values, I have to make assumptions. The items available are 3 apples, 1 soap, 1 egg. The other player could value any of these. If I take the egg (which I value highest), that leaves 3 apples and 1 soap. The other player might want the soap and some apples. If they take the soap and all 3 apples, that's possible. But if they also want the egg, there's a conflict. However, since I can't know their preferences, I need to structure my proposal in a way that allows the other player to take what's left without overlapping.} \\
\\ 
\texttt{Alternatively, maybe I should take the soap and some apples. If I take the soap (3 points) and, say, 1 apple (1 point), that's 4 points total. Then the other player could take the remaining 2 apples and the egg. But if the other player values the egg highly, they might want it, which would conflict if I take the soap. Wait, but if I take the soap and 1 apple, and they take the egg and 2 apples, that's possible. But my total would be 3+1=4, which is better than just taking the egg (4) alone. Hmm, same points. But maybe taking the egg is better because it's a single item, leaving more flexibility for the other player.} \\
\\ 
\texttt{Another angle: if I take the egg, the other player can take the soap and all 3 apples. That would give me 4 points, and the other player would get their value for soap and apples. If their value for soap is high, that's good for them. If they value apples more, maybe they prefer more apples. Either way, as long as their proposal doesn't require the egg, it's okay.} \\
\\ 
\texttt{But what if the other player also wants the egg? Then our proposals would conflict. But since I don't know their values, I have to assume they might not. Maybe the best strategy is to take the highest value item for} \\
            }
        }
    }
     \\ \\

    \theutterance \stepcounter{utterance}  
    & & & \multicolumn{2}{p{0.3\linewidth}}{
        \cellcolor[rgb]{0.9,0.9,0.9}{
            \makecell[{{p{\linewidth}}}]{
                \texttt{\tiny{[GM$|$GM]}}
                \texttt{continue} \\
            }
        }
    }
    & & \\ \\

    \theutterance \stepcounter{utterance}  
    & & & \multicolumn{2}{p{0.3\linewidth}}{
        \cellcolor[rgb]{0.9,0.9,0.9}{
            \makecell[{{p{\linewidth}}}]{
                \texttt{\tiny{[GM$|$GM]}}
                \texttt{<think>} \\
\texttt{Okay, let's try to figure out the best way to approach this. My values are apple:1, soap:3, egg:4. The other player's values are unknown. The key is to make a proposal that maximizes my points while ensuring that the other player's proposal doesn't conflict with mine.} \\
\\ 
\texttt{First, I need to consider what the other player might want. Since I don't know their values, I have to make assumptions. The items available are 3 apples, 1 soap, 1 egg. The other player could value any of these. If I take the egg (which I value highest), that leaves 3 apples and 1 soap. The other player might want the soap and some apples. If they take the soap and all 3 apples, that's possible. But if they also want the egg, there's a conflict. However, since I can't know their preferences, I need to structure my proposal in a way that allows the other player to take what's left without overlapping.} \\
\\ 
\texttt{Alternatively, maybe I should take the soap and some apples. If I take the soap (3 points) and, say, 1 apple (1 point), that's 4 points total. Then the other player could take the remaining 2 apples and the egg. But if the other player values the egg highly, they might want it, which would conflict if I take the soap. Wait, but if I take the soap and 1 apple, and they take the egg and 2 apples, that's possible. But my total would be 3+1=4, which is better than just taking the egg (4) alone. Hmm, same points. But maybe taking the egg is better because it's a single item, leaving more flexibility for the other player.} \\
\\ 
\texttt{Another angle: if I take the egg, the other player can take the soap and all 3 apples. That would give me 4 points, and the other player would get their value for soap and apples. If their value for soap is high, that's good for them. If they value apples more, maybe they prefer more apples. Either way, as long as their proposal doesn't require the egg, it's okay.} \\
\\ 
\texttt{But what if the other player also wants the egg? Then our proposals would conflict. But since I don't know their values, I have to assume they might not. Maybe the best strategy is to take the highest value item for} \\
            }
        }
    }
    & & \\ \\

    \theutterance \stepcounter{utterance}  
    & & \multicolumn{4}{p{0.6\linewidth}}{
        \cellcolor[rgb]{0.9,0.9,0.9}{
            \makecell[{{p{\linewidth}}}]{
                \texttt{\tiny{[P1$\langle$GM]}}
                \texttt{<think>} \\
\texttt{Okay, let's try to figure out the best way to approach this. My values are apple:1, soap:3, egg:4. The other player's values are unknown. The key is to make a proposal that maximizes my points while ensuring that the other player's proposal doesn't conflict with mine.} \\
\\ 
\texttt{First, I need to consider what the other player might want. Since I don't know their values, I have to make assumptions. The items available are 3 apples, 1 soap, 1 egg. The other player could value any of these. If I take the egg (which I value highest), that leaves 3 apples and 1 soap. The other player might want the soap and some apples. If they take the soap and all 3 apples, that's possible. But if they also want the egg, there's a conflict. However, since I can't know their preferences, I need to structure my proposal in a way that allows the other player to take what's left without overlapping.} \\
\\ 
\texttt{Alternatively, maybe I should take the soap and some apples. If I take the soap (3 points) and, say, 1 apple (1 point), that's 4 points total. Then the other player could take the remaining 2 apples and the egg. But if the other player values the egg highly, they might want it, which would conflict if I take the soap. Wait, but if I take the soap and 1 apple, and they take the egg and 2 apples, that's possible. But my total would be 3+1=4, which is better than just taking the egg (4) alone. Hmm, same points. But maybe taking the egg is better because it's a single item, leaving more flexibility for the other player.} \\
\\ 
\texttt{Another angle: if I take the egg, the other player can take the soap and all 3 apples. That would give me 4 points, and the other player would get their value for soap and apples. If their value for soap is high, that's good for them. If they value apples more, maybe they prefer more apples. Either way, as long as their proposal doesn't require the egg, it's okay.} \\
\\ 
\texttt{But what if the other player also wants the egg? Then our proposals would conflict. But since I don't know their values, I have to assume they might not. Maybe the best strategy is to take the highest value item for} \\
            }
        }
    }
    & \\ \\

    \theutterance \stepcounter{utterance}  
    & \multicolumn{4}{p{0.6\linewidth}}{
        \cellcolor[rgb]{0.9,0.9,0.9}{
            \makecell[{{p{\linewidth}}}]{
                \texttt{\tiny{[P1$\rangle$GM]}}
                \texttt{<think>} \\
\texttt{Okay, let's try to figure out the best way to approach this. My values are apple:1, soap:3, egg:4. The other player's values are unknown. The key is to make a proposal that maximizes my points while ensuring that the other player's proposal doesn't conflict with mine.} \\
\\ 
\texttt{First, I need to consider what the other player might want. Since I don't know their values, I have to make assumptions. The items available are 3 apples, 1 soap, 1 egg. The other player could value any of these. If I take the egg (which I value highest), that leaves 3 apples and 1 soap. The other player might want the soap and some apples. If they take the soap and all 3 apples, that's possible. But if they also want the egg, there's a conflict. However, since I can't know their preferences, I need to structure my proposal in a way that allows the other player to take what's left without overlapping.} \\
\\ 
\texttt{Alternatively, maybe I should take the soap and some apples. If I take the soap (3 points) and, say, 1 apple (1 point), that's 4 points total. Then the other player could take the remaining 2 apples and the egg. But if the other player values the egg highly, they might want it, which would conflict if I take the soap. Wait, but if I take the soap and 1 apple, and they take the egg and 2 apples, that's possible. But my total would be 3+1=4, which is better than just taking the egg (4) alone. Hmm, same points. But maybe taking the egg is better because it's a single item, leaving more flexibility for the other player.} \\
\\ 
\texttt{Another angle: if I take the egg, the other player can take the soap and all 3 apples. That would give me 4 points, and the other player would get their value for soap and apples. If their value for soap is high, that's good for them. If they value apples more, maybe they prefer more apples. Either way, as long as their proposal doesn't require the egg, it's okay.} \\
\\ 
\texttt{But what if the other player also wants the egg? Then our proposals would conflict. But since I don't know their values, I have to assume they might not. Maybe the best strategy is to take the highest value item for} \\
            }
        }
    }
    & & \\ \\

    \theutterance \stepcounter{utterance}  
    & & & \multicolumn{2}{p{0.3\linewidth}}{
        \cellcolor[rgb]{0.9,0.9,0.9}{
            \makecell[{{p{\linewidth}}}]{
                \texttt{\tiny{[GM$|$GM]}}
                \texttt{continue} \\
            }
        }
    }
    & & \\ \\

    \theutterance \stepcounter{utterance}  
    & & & \multicolumn{2}{p{0.3\linewidth}}{
        \cellcolor[rgb]{0.9,0.9,0.9}{
            \makecell[{{p{\linewidth}}}]{
                \texttt{\tiny{[GM$|$GM]}}
                \texttt{<think>} \\
\texttt{Okay, let's try to figure out the best way to approach this. My values are apple:1, soap:3, egg:4. The other player's values are unknown. The key is to make a proposal that maximizes my points while ensuring that the other player's proposal doesn't conflict with mine.} \\
\\ 
\texttt{First, I need to consider what the other player might want. Since I don't know their values, I have to make assumptions. The items available are 3 apples, 1 soap, 1 egg. The other player could value any of these. If I take the egg (which I value highest), that leaves 3 apples and 1 soap. The other player might want the soap and some apples. If they take the soap and all 3 apples, that's possible. But if they also want the egg, there's a conflict. However, since I can't know their preferences, I need to structure my proposal in a way that allows the other player to take what's left without overlapping.} \\
\\ 
\texttt{Alternatively, maybe I should take the soap and some apples. If I take the soap (3 points) and, say, 1 apple (1 point), that's 4 points total. Then the other player could take the remaining 2 apples and the egg. But if the other player values the egg highly, they might want it, which would conflict if I take the soap. Wait, but if I take the soap and 1 apple, and they take the egg and 2 apples, that's possible. But my total would be 3+1=4, which is better than just taking the egg (4) alone. Hmm, same points. But maybe taking the egg is better because it's a single item, leaving more flexibility for the other player.} \\
\\ 
\texttt{Another angle: if I take the egg, the other player can take the soap and all 3 apples. That would give me 4 points, and the other player would get their value for soap and apples. If their value for soap is high, that's good for them. If they value apples more, maybe they prefer more apples. Either way, as long as their proposal doesn't require the egg, it's okay.} \\
\\ 
\texttt{But what if the other player also wants the egg? Then our proposals would conflict. But since I don't know their values, I have to assume they might not. Maybe the best strategy is to take the highest value item for} \\
            }
        }
    }
    & & \\ \\

    \theutterance \stepcounter{utterance}  
    & & \multicolumn{4}{p{0.6\linewidth}}{
        \cellcolor[rgb]{0.9,0.9,0.9}{
            \makecell[{{p{\linewidth}}}]{
                \texttt{\tiny{[GM$\rangle$P2]}}
                \texttt{<think>} \\
\texttt{Okay, let's try to figure out the best way to approach this. My values are apple:1, soap:3, egg:4. The other player's values are unknown. The key is to make a proposal that maximizes my points while ensuring that the other player's proposal doesn't conflict with mine.} \\
\\ 
\texttt{First, I need to consider what the other player might want. Since I don't know their values, I have to make assumptions. The items available are 3 apples, 1 soap, 1 egg. The other player could value any of these. If I take the egg (which I value highest), that leaves 3 apples and 1 soap. The other player might want the soap and some apples. If they take the soap and all 3 apples, that's possible. But if they also want the egg, there's a conflict. However, since I can't know their preferences, I need to structure my proposal in a way that allows the other player to take what's left without overlapping.} \\
\\ 
\texttt{Alternatively, maybe I should take the soap and some apples. If I take the soap (3 points) and, say, 1 apple (1 point), that's 4 points total. Then the other player could take the remaining 2 apples and the egg. But if the other player values the egg highly, they might want it, which would conflict if I take the soap. Wait, but if I take the soap and 1 apple, and they take the egg and 2 apples, that's possible. But my total would be 3+1=4, which is better than just taking the egg (4) alone. Hmm, same points. But maybe taking the egg is better because it's a single item, leaving more flexibility for the other player.} \\
\\ 
\texttt{Another angle: if I take the egg, the other player can take the soap and all 3 apples. That would give me 4 points, and the other player would get their value for soap and apples. If their value for soap is high, that's good for them. If they value apples more, maybe they prefer more apples. Either way, as long as their proposal doesn't require the egg, it's okay.} \\
\\ 
\texttt{But what if the other player also wants the egg? Then our proposals would conflict. But since I don't know their values, I have to assume they might not. Maybe the best strategy is to take the highest value item for} \\
            }
        }
    }
    & \\ \\

    \theutterance \stepcounter{utterance}  
    & & & \multicolumn{4}{p{0.6\linewidth}}{
        \cellcolor[rgb]{0.9,0.9,0.9}{
            \makecell[{{p{\linewidth}}}]{
                \texttt{\tiny{[GM$\langle$P2]}}
                \texttt{<think>} \\
\texttt{Okay, let's try to figure out the best way to approach this. My values are apple:1, soap:3, egg:4. The other player's values are unknown. The key is to make a proposal that maximizes my points while ensuring that the other player's proposal doesn't conflict with mine.} \\
\\ 
\texttt{First, I need to consider what the other player might want. Since I don't know their values, I have to make assumptions. The items available are 3 apples, 1 soap, 1 egg. The other player could value any of these. If I take the egg (which I value highest), that leaves 3 apples and 1 soap. The other player might want the soap and some apples. If they take the soap and all 3 apples, that's possible. But if they also want the egg, there's a conflict. However, since I can't know their preferences, I need to structure my proposal in a way that allows the other player to take what's left without overlapping.} \\
\\ 
\texttt{Alternatively, maybe I should take the soap and some apples. If I take the soap (3 points) and, say, 1 apple (1 point), that's 4 points total. Then the other player could take the remaining 2 apples and the egg. But if the other player values the egg highly, they might want it, which would conflict if I take the soap. Wait, but if I take the soap and 1 apple, and they take the egg and 2 apples, that's possible. But my total would be 3+1=4, which is better than just taking the egg (4) alone. Hmm, same points. But maybe taking the egg is better because it's a single item, leaving more flexibility for the other player.} \\
\\ 
\texttt{Another angle: if I take the egg, the other player can take the soap and all 3 apples. That would give me 4 points, and the other player would get their value for soap and apples. If their value for soap is high, that's good for them. If they value apples more, maybe they prefer more apples. Either way, as long as their proposal doesn't require the egg, it's okay.} \\
\\ 
\texttt{But what if the other player also wants the egg? Then our proposals would conflict. But since I don't know their values, I have to assume they might not. Maybe the best strategy is to take the highest value item for} \\
            }
        }
    }
     \\ \\

    \theutterance \stepcounter{utterance}  
    & & & \multicolumn{2}{p{0.3\linewidth}}{
        \cellcolor[rgb]{0.9,0.9,0.9}{
            \makecell[{{p{\linewidth}}}]{
                \texttt{\tiny{[GM$|$GM]}}
                \texttt{continue} \\
            }
        }
    }
    & & \\ \\

    \theutterance \stepcounter{utterance}  
    & & & \multicolumn{2}{p{0.3\linewidth}}{
        \cellcolor[rgb]{0.9,0.9,0.9}{
            \makecell[{{p{\linewidth}}}]{
                \texttt{\tiny{[GM$|$GM]}}
                \texttt{<think>} \\
\texttt{Okay, let's try to figure out the best way to approach this. My values are apple:1, soap:3, egg:4. The other player's values are unknown. The key is to make a proposal that maximizes my points while ensuring that the other player's proposal doesn't conflict with mine.} \\
\\ 
\texttt{First, I need to consider what the other player might want. Since I don't know their values, I have to make assumptions. The items available are 3 apples, 1 soap, 1 egg. The other player could value any of these. If I take the egg (which I value highest), that leaves 3 apples and 1 soap. The other player might want the soap and some apples. If they take the soap and all 3 apples, that's possible. But if they also want the egg, there's a conflict. However, since I can't know their preferences, I need to structure my proposal in a way that allows the other player to take what's left without overlapping.} \\
\\ 
\texttt{Alternatively, maybe I should take the soap and some apples. If I take the soap (3 points) and, say, 1 apple (1 point), that's 4 points total. Then the other player could take the remaining 2 apples and the egg. But if the other player values the egg highly, they might want it, which would conflict if I take the soap. Wait, but if I take the soap and 1 apple, and they take the egg and 2 apples, that's possible. But my total would be 3+1=4, which is better than just taking the egg (4) alone. Hmm, same points. But maybe taking the egg is better because it's a single item, leaving more flexibility for the other player.} \\
\\ 
\texttt{Another angle: if I take the egg, the other player can take the soap and all 3 apples. That would give me 4 points, and the other player would get their value for soap and apples. If their value for soap is high, that's good for them. If they value apples more, maybe they prefer more apples. Either way, as long as their proposal doesn't require the egg, it's okay.} \\
\\ 
\texttt{But what if the other player also wants the egg? Then our proposals would conflict. But since I don't know their values, I have to assume they might not. Maybe the best strategy is to take the highest value item for} \\
            }
        }
    }
    & & \\ \\

    \theutterance \stepcounter{utterance}  
    & & \multicolumn{4}{p{0.6\linewidth}}{
        \cellcolor[rgb]{0.9,0.9,0.9}{
            \makecell[{{p{\linewidth}}}]{
                \texttt{\tiny{[P1$\langle$GM]}}
                \texttt{<think>} \\
\texttt{Okay, let's try to figure out the best way to approach this. My values are apple:1, soap:3, egg:4. The other player's values are unknown. The key is to make a proposal that maximizes my points while ensuring that the other player's proposal doesn't conflict with mine.} \\
\\ 
\texttt{First, I need to consider what the other player might want. Since I don't know their values, I have to make assumptions. The items available are 3 apples, 1 soap, 1 egg. The other player could value any of these. If I take the egg (which I value highest), that leaves 3 apples and 1 soap. The other player might want the soap and some apples. If they take the soap and all 3 apples, that's possible. But if they also want the egg, there's a conflict. However, since I can't know their preferences, I need to structure my proposal in a way that allows the other player to take what's left without overlapping.} \\
\\ 
\texttt{Alternatively, maybe I should take the soap and some apples. If I take the soap (3 points) and, say, 1 apple (1 point), that's 4 points total. Then the other player could take the remaining 2 apples and the egg. But if the other player values the egg highly, they might want it, which would conflict if I take the soap. Wait, but if I take the soap and 1 apple, and they take the egg and 2 apples, that's possible. But my total would be 3+1=4, which is better than just taking the egg (4) alone. Hmm, same points. But maybe taking the egg is better because it's a single item, leaving more flexibility for the other player.} \\
\\ 
\texttt{Another angle: if I take the egg, the other player can take the soap and all 3 apples. That would give me 4 points, and the other player would get their value for soap and apples. If their value for soap is high, that's good for them. If they value apples more, maybe they prefer more apples. Either way, as long as their proposal doesn't require the egg, it's okay.} \\
\\ 
\texttt{But what if the other player also wants the egg? Then our proposals would conflict. But since I don't know their values, I have to assume they might not. Maybe the best strategy is to take the highest value item for} \\
            }
        }
    }
    & \\ \\

    \theutterance \stepcounter{utterance}  
    & \multicolumn{4}{p{0.6\linewidth}}{
        \cellcolor[rgb]{0.9,0.9,0.9}{
            \makecell[{{p{\linewidth}}}]{
                \texttt{\tiny{[P1$\rangle$GM]}}
                \texttt{<think>} \\
\texttt{Okay, let's try to figure out the best way to approach this. My values are apple:1, soap:3, egg:4. The other player's values are unknown. The key is to make a proposal that maximizes my points while ensuring that the other player's proposal doesn't conflict with mine.} \\
\\ 
\texttt{First, I need to consider what the other player might want. Since I don't know their values, I have to make assumptions. The items available are 3 apples, 1 soap, 1 egg. The other player could value any of these. If I take the egg (which I value highest), that leaves 3 apples and 1 soap. The other player might want the soap and some apples. If they take the soap and all 3 apples, that's possible. But if they also want the egg, there's a conflict. However, since I can't know their preferences, I need to structure my proposal in a way that allows the other player to take what's left without overlapping.} \\
\\ 
\texttt{Alternatively, maybe I should take the soap and some apples. If I take the soap (3 points) and, say, 1 apple (1 point), that's 4 points total. Then the other player could take the remaining 2 apples and the egg. But if the other player values the egg highly, they might want it, which would conflict if I take the soap. Wait, but if I take the soap and 1 apple, and they take the egg and 2 apples, that's possible. But my total would be 3+1=4, which is better than just taking the egg (4) alone. Hmm, same points. But maybe taking the egg is better because it's a single item, leaving more flexibility for the other player.} \\
\\ 
\texttt{Another angle: if I take the egg, the other player can take the soap and all 3 apples. That would give me 4 points, and the other player would get their value for soap and apples. If their value for soap is high, that's good for them. If they value apples more, maybe they prefer more apples. Either way, as long as their proposal doesn't require the egg, it's okay.} \\
\\ 
\texttt{But what if the other player also wants the egg? Then our proposals would conflict. But since I don't know their values, I have to assume they might not. Maybe the best strategy is to take the highest value item for} \\
            }
        }
    }
    & & \\ \\

    \theutterance \stepcounter{utterance}  
    & & & \multicolumn{2}{p{0.3\linewidth}}{
        \cellcolor[rgb]{0.9,0.9,0.9}{
            \makecell[{{p{\linewidth}}}]{
                \texttt{\tiny{[GM$|$GM]}}
                \texttt{continue} \\
            }
        }
    }
    & & \\ \\

    \theutterance \stepcounter{utterance}  
    & & & \multicolumn{2}{p{0.3\linewidth}}{
        \cellcolor[rgb]{0.9,0.9,0.9}{
            \makecell[{{p{\linewidth}}}]{
                \texttt{\tiny{[GM$|$GM]}}
                \texttt{<think>} \\
\texttt{Okay, let's try to figure out the best way to approach this. My values are apple:1, soap:3, egg:4. The other player's values are unknown. The key is to make a proposal that maximizes my points while ensuring that the other player's proposal doesn't conflict with mine.} \\
\\ 
\texttt{First, I need to consider what the other player might want. Since I don't know their values, I have to make assumptions. The items available are 3 apples, 1 soap, 1 egg. The other player could value any of these. If I take the egg (which I value highest), that leaves 3 apples and 1 soap. The other player might want the soap and some apples. If they take the soap and all 3 apples, that's possible. But if they also want the egg, there's a conflict. However, since I can't know their preferences, I need to structure my proposal in a way that allows the other player to take what's left without overlapping.} \\
\\ 
\texttt{Alternatively, maybe I should take the soap and some apples. If I take the soap (3 points) and, say, 1 apple (1 point), that's 4 points total. Then the other player could take the remaining 2 apples and the egg. But if the other player values the egg highly, they might want it, which would conflict if I take the soap. Wait, but if I take the soap and 1 apple, and they take the egg and 2 apples, that's possible. But my total would be 3+1=4, which is better than just taking the egg (4) alone. Hmm, same points. But maybe taking the egg is better because it's a single item, leaving more flexibility for the other player.} \\
\\ 
\texttt{Another angle: if I take the egg, the other player can take the soap and all 3 apples. That would give me 4 points, and the other player would get their value for soap and apples. If their value for soap is high, that's good for them. If they value apples more, maybe they prefer more apples. Either way, as long as their proposal doesn't require the egg, it's okay.} \\
\\ 
\texttt{But what if the other player also wants the egg? Then our proposals would conflict. But since I don't know their values, I have to assume they might not. Maybe the best strategy is to take the highest value item for} \\
            }
        }
    }
    & & \\ \\

    \theutterance \stepcounter{utterance}  
    & & \multicolumn{4}{p{0.6\linewidth}}{
        \cellcolor[rgb]{0.9,0.9,0.9}{
            \makecell[{{p{\linewidth}}}]{
                \texttt{\tiny{[GM$\rangle$P2]}}
                \texttt{<think>} \\
\texttt{Okay, let's try to figure out the best way to approach this. My values are apple:1, soap:3, egg:4. The other player's values are unknown. The key is to make a proposal that maximizes my points while ensuring that the other player's proposal doesn't conflict with mine.} \\
\\ 
\texttt{First, I need to consider what the other player might want. Since I don't know their values, I have to make assumptions. The items available are 3 apples, 1 soap, 1 egg. The other player could value any of these. If I take the egg (which I value highest), that leaves 3 apples and 1 soap. The other player might want the soap and some apples. If they take the soap and all 3 apples, that's possible. But if they also want the egg, there's a conflict. However, since I can't know their preferences, I need to structure my proposal in a way that allows the other player to take what's left without overlapping.} \\
\\ 
\texttt{Alternatively, maybe I should take the soap and some apples. If I take the soap (3 points) and, say, 1 apple (1 point), that's 4 points total. Then the other player could take the remaining 2 apples and the egg. But if the other player values the egg highly, they might want it, which would conflict if I take the soap. Wait, but if I take the soap and 1 apple, and they take the egg and 2 apples, that's possible. But my total would be 3+1=4, which is better than just taking the egg (4) alone. Hmm, same points. But maybe taking the egg is better because it's a single item, leaving more flexibility for the other player.} \\
\\ 
\texttt{Another angle: if I take the egg, the other player can take the soap and all 3 apples. That would give me 4 points, and the other player would get their value for soap and apples. If their value for soap is high, that's good for them. If they value apples more, maybe they prefer more apples. Either way, as long as their proposal doesn't require the egg, it's okay.} \\
\\ 
\texttt{But what if the other player also wants the egg? Then our proposals would conflict. But since I don't know their values, I have to assume they might not. Maybe the best strategy is to take the highest value item for} \\
            }
        }
    }
    & \\ \\

    \theutterance \stepcounter{utterance}  
    & & & \multicolumn{4}{p{0.6\linewidth}}{
        \cellcolor[rgb]{0.9,0.9,0.9}{
            \makecell[{{p{\linewidth}}}]{
                \texttt{\tiny{[GM$\langle$P2]}}
                \texttt{<think>} \\
\texttt{Okay, let's try to figure out the best way to approach this. My values are apple:1, soap:3, egg:4. The other player's values are unknown. The key is to make a proposal that maximizes my points while ensuring that the other player's proposal doesn't conflict with mine.} \\
\\ 
\texttt{First, I need to consider what the other player might want. Since I don't know their values, I have to make assumptions. The items available are 3 apples, 1 soap, 1 egg. The other player could value any of these. If I take the egg (which I value highest), that leaves 3 apples and 1 soap. The other player might want the soap and some apples. If they take the soap and all 3 apples, that's possible. But if they also want the egg, there's a conflict. However, since I can't know their preferences, I need to structure my proposal in a way that allows the other player to take what's left without overlapping.} \\
\\ 
\texttt{Alternatively, maybe I should take the soap and some apples. If I take the soap (3 points) and, say, 1 apple (1 point), that's 4 points total. Then the other player could take the remaining 2 apples and the egg. But if the other player values the egg highly, they might want it, which would conflict if I take the soap. Wait, but if I take the soap and 1 apple, and they take the egg and 2 apples, that's possible. But my total would be 3+1=4, which is better than just taking the egg (4) alone. Hmm, same points. But maybe taking the egg is better because it's a single item, leaving more flexibility for the other player.} \\
\\ 
\texttt{Another angle: if I take the egg, the other player can take the soap and all 3 apples. That would give me 4 points, and the other player would get their value for soap and apples. If their value for soap is high, that's good for them. If they value apples more, maybe they prefer more apples. Either way, as long as their proposal doesn't require the egg, it's okay.} \\
\\ 
\texttt{But what if the other player also wants the egg? Then our proposals would conflict. But since I don't know their values, I have to assume they might not. Maybe the best strategy is to take the highest value item for} \\
            }
        }
    }
     \\ \\

    \theutterance \stepcounter{utterance}  
    & & & \multicolumn{2}{p{0.3\linewidth}}{
        \cellcolor[rgb]{0.9,0.9,0.9}{
            \makecell[{{p{\linewidth}}}]{
                \texttt{\tiny{[GM$|$GM]}}
                \texttt{continue} \\
            }
        }
    }
    & & \\ \\

    \theutterance \stepcounter{utterance}  
    & & & \multicolumn{2}{p{0.3\linewidth}}{
        \cellcolor[rgb]{0.9,0.9,0.9}{
            \makecell[{{p{\linewidth}}}]{
                \texttt{\tiny{[GM$|$GM]}}
                \texttt{<think>} \\
\texttt{Okay, let's try to figure out the best way to approach this. My values are apple:1, soap:3, egg:4. The other player's values are unknown. The key is to make a proposal that maximizes my points while ensuring that the other player's proposal doesn't conflict with mine.} \\
\\ 
\texttt{First, I need to consider what the other player might want. Since I don't know their values, I have to make assumptions. The items available are 3 apples, 1 soap, 1 egg. The other player could value any of these. If I take the egg (which I value highest), that leaves 3 apples and 1 soap. The other player might want the soap and some apples. If they take the soap and all 3 apples, that's possible. But if they also want the egg, there's a conflict. However, since I can't know their preferences, I need to structure my proposal in a way that allows the other player to take what's left without overlapping.} \\
\\ 
\texttt{Alternatively, maybe I should take the soap and some apples. If I take the soap (3 points) and, say, 1 apple (1 point), that's 4 points total. Then the other player could take the remaining 2 apples and the egg. But if the other player values the egg highly, they might want it, which would conflict if I take the soap. Wait, but if I take the soap and 1 apple, and they take the egg and 2 apples, that's possible. But my total would be 3+1=4, which is better than just taking the egg (4) alone. Hmm, same points. But maybe taking the egg is better because it's a single item, leaving more flexibility for the other player.} \\
\\ 
\texttt{Another angle: if I take the egg, the other player can take the soap and all 3 apples. That would give me 4 points, and the other player would get their value for soap and apples. If their value for soap is high, that's good for them. If they value apples more, maybe they prefer more apples. Either way, as long as their proposal doesn't require the egg, it's okay.} \\
\\ 
\texttt{But what if the other player also wants the egg? Then our proposals would conflict. But since I don't know their values, I have to assume they might not. Maybe the best strategy is to take the highest value item for} \\
            }
        }
    }
    & & \\ \\

    \theutterance \stepcounter{utterance}  
    & & \multicolumn{4}{p{0.6\linewidth}}{
        \cellcolor[rgb]{0.9,0.9,0.9}{
            \makecell[{{p{\linewidth}}}]{
                \texttt{\tiny{[P1$\langle$GM]}}
                \texttt{<think>} \\
\texttt{Okay, let's try to figure out the best way to approach this. My values are apple:1, soap:3, egg:4. The other player's values are unknown. The key is to make a proposal that maximizes my points while ensuring that the other player's proposal doesn't conflict with mine.} \\
\\ 
\texttt{First, I need to consider what the other player might want. Since I don't know their values, I have to make assumptions. The items available are 3 apples, 1 soap, 1 egg. The other player could value any of these. If I take the egg (which I value highest), that leaves 3 apples and 1 soap. The other player might want the soap and some apples. If they take the soap and all 3 apples, that's possible. But if they also want the egg, there's a conflict. However, since I can't know their preferences, I need to structure my proposal in a way that allows the other player to take what's left without overlapping.} \\
\\ 
\texttt{Alternatively, maybe I should take the soap and some apples. If I take the soap (3 points) and, say, 1 apple (1 point), that's 4 points total. Then the other player could take the remaining 2 apples and the egg. But if the other player values the egg highly, they might want it, which would conflict if I take the soap. Wait, but if I take the soap and 1 apple, and they take the egg and 2 apples, that's possible. But my total would be 3+1=4, which is better than just taking the egg (4) alone. Hmm, same points. But maybe taking the egg is better because it's a single item, leaving more flexibility for the other player.} \\
\\ 
\texttt{Another angle: if I take the egg, the other player can take the soap and all 3 apples. That would give me 4 points, and the other player would get their value for soap and apples. If their value for soap is high, that's good for them. If they value apples more, maybe they prefer more apples. Either way, as long as their proposal doesn't require the egg, it's okay.} \\
\\ 
\texttt{But what if the other player also wants the egg? Then our proposals would conflict. But since I don't know their values, I have to assume they might not. Maybe the best strategy is to take the highest value item for} \\
            }
        }
    }
    & \\ \\

    \theutterance \stepcounter{utterance}  
    & \multicolumn{4}{p{0.6\linewidth}}{
        \cellcolor[rgb]{0.9,0.9,0.9}{
            \makecell[{{p{\linewidth}}}]{
                \texttt{\tiny{[P1$\rangle$GM]}}
                \texttt{<think>} \\
\texttt{Okay, let's try to figure out the best way to approach this. My values are apple:1, soap:3, egg:4. The other player's values are unknown. The key is to make a proposal that maximizes my points while ensuring that the other player's proposal doesn't conflict with mine.} \\
\\ 
\texttt{First, I need to consider what the other player might want. Since I don't know their values, I have to make assumptions. The items available are 3 apples, 1 soap, 1 egg. The other player could value any of these. If I take the egg (which I value highest), that leaves 3 apples and 1 soap. The other player might want the soap and some apples. If they take the soap and all 3 apples, that's possible. But if they also want the egg, there's a conflict. However, since I can't know their preferences, I need to structure my proposal in a way that allows the other player to take what's left without overlapping.} \\
\\ 
\texttt{Alternatively, maybe I should take the soap and some apples. If I take the soap (3 points) and, say, 1 apple (1 point), that's 4 points total. Then the other player could take the remaining 2 apples and the egg. But if the other player values the egg highly, they might want it, which would conflict if I take the soap. Wait, but if I take the soap and 1 apple, and they take the egg and 2 apples, that's possible. But my total would be 3+1=4, which is better than just taking the egg (4) alone. Hmm, same points. But maybe taking the egg is better because it's a single item, leaving more flexibility for the other player.} \\
\\ 
\texttt{Another angle: if I take the egg, the other player can take the soap and all 3 apples. That would give me 4 points, and the other player would get their value for soap and apples. If their value for soap is high, that's good for them. If they value apples more, maybe they prefer more apples. Either way, as long as their proposal doesn't require the egg, it's okay.} \\
\\ 
\texttt{But what if the other player also wants the egg? Then our proposals would conflict. But since I don't know their values, I have to assume they might not. Maybe the best strategy is to take the highest value item for} \\
            }
        }
    }
    & & \\ \\

    \theutterance \stepcounter{utterance}  
    & & & \multicolumn{2}{p{0.3\linewidth}}{
        \cellcolor[rgb]{0.9,0.9,0.9}{
            \makecell[{{p{\linewidth}}}]{
                \texttt{\tiny{[GM$|$GM]}}
                \texttt{continue} \\
            }
        }
    }
    & & \\ \\

    \theutterance \stepcounter{utterance}  
    & & & \multicolumn{2}{p{0.3\linewidth}}{
        \cellcolor[rgb]{0.9,0.9,0.9}{
            \makecell[{{p{\linewidth}}}]{
                \texttt{\tiny{[GM$|$GM]}}
                \texttt{<think>} \\
\texttt{Okay, let's try to figure out the best way to approach this. My values are apple:1, soap:3, egg:4. The other player's values are unknown. The key is to make a proposal that maximizes my points while ensuring that the other player's proposal doesn't conflict with mine.} \\
\\ 
\texttt{First, I need to consider what the other player might want. Since I don't know their values, I have to make assumptions. The items available are 3 apples, 1 soap, 1 egg. The other player could value any of these. If I take the egg (which I value highest), that leaves 3 apples and 1 soap. The other player might want the soap and some apples. If they take the soap and all 3 apples, that's possible. But if they also want the egg, there's a conflict. However, since I can't know their preferences, I need to structure my proposal in a way that allows the other player to take what's left without overlapping.} \\
\\ 
\texttt{Alternatively, maybe I should take the soap and some apples. If I take the soap (3 points) and, say, 1 apple (1 point), that's 4 points total. Then the other player could take the remaining 2 apples and the egg. But if the other player values the egg highly, they might want it, which would conflict if I take the soap. Wait, but if I take the soap and 1 apple, and they take the egg and 2 apples, that's possible. But my total would be 3+1=4, which is better than just taking the egg (4) alone. Hmm, same points. But maybe taking the egg is better because it's a single item, leaving more flexibility for the other player.} \\
\\ 
\texttt{Another angle: if I take the egg, the other player can take the soap and all 3 apples. That would give me 4 points, and the other player would get their value for soap and apples. If their value for soap is high, that's good for them. If they value apples more, maybe they prefer more apples. Either way, as long as their proposal doesn't require the egg, it's okay.} \\
\\ 
\texttt{But what if the other player also wants the egg? Then our proposals would conflict. But since I don't know their values, I have to assume they might not. Maybe the best strategy is to take the highest value item for} \\
            }
        }
    }
    & & \\ \\

    \theutterance \stepcounter{utterance}  
    & & \multicolumn{4}{p{0.6\linewidth}}{
        \cellcolor[rgb]{0.9,0.9,0.9}{
            \makecell[{{p{\linewidth}}}]{
                \texttt{\tiny{[GM$\rangle$P2]}}
                \texttt{<think>} \\
\texttt{Okay, let's try to figure out the best way to approach this. My values are apple:1, soap:3, egg:4. The other player's values are unknown. The key is to make a proposal that maximizes my points while ensuring that the other player's proposal doesn't conflict with mine.} \\
\\ 
\texttt{First, I need to consider what the other player might want. Since I don't know their values, I have to make assumptions. The items available are 3 apples, 1 soap, 1 egg. The other player could value any of these. If I take the egg (which I value highest), that leaves 3 apples and 1 soap. The other player might want the soap and some apples. If they take the soap and all 3 apples, that's possible. But if they also want the egg, there's a conflict. However, since I can't know their preferences, I need to structure my proposal in a way that allows the other player to take what's left without overlapping.} \\
\\ 
\texttt{Alternatively, maybe I should take the soap and some apples. If I take the soap (3 points) and, say, 1 apple (1 point), that's 4 points total. Then the other player could take the remaining 2 apples and the egg. But if the other player values the egg highly, they might want it, which would conflict if I take the soap. Wait, but if I take the soap and 1 apple, and they take the egg and 2 apples, that's possible. But my total would be 3+1=4, which is better than just taking the egg (4) alone. Hmm, same points. But maybe taking the egg is better because it's a single item, leaving more flexibility for the other player.} \\
\\ 
\texttt{Another angle: if I take the egg, the other player can take the soap and all 3 apples. That would give me 4 points, and the other player would get their value for soap and apples. If their value for soap is high, that's good for them. If they value apples more, maybe they prefer more apples. Either way, as long as their proposal doesn't require the egg, it's okay.} \\
\\ 
\texttt{But what if the other player also wants the egg? Then our proposals would conflict. But since I don't know their values, I have to assume they might not. Maybe the best strategy is to take the highest value item for} \\
            }
        }
    }
    & \\ \\

    \theutterance \stepcounter{utterance}  
    & & & \multicolumn{4}{p{0.6\linewidth}}{
        \cellcolor[rgb]{0.9,0.9,0.9}{
            \makecell[{{p{\linewidth}}}]{
                \texttt{\tiny{[GM$\langle$P2]}}
                \texttt{<think>} \\
\texttt{Okay, let's try to figure out the best way to approach this. My values are apple:1, soap:3, egg:4. The other player's values are unknown. The key is to make a proposal that maximizes my points while ensuring that the other player's proposal doesn't conflict with mine.} \\
\\ 
\texttt{First, I need to consider what the other player might want. Since I don't know their values, I have to make assumptions. The items available are 3 apples, 1 soap, 1 egg. The other player could value any of these. If I take the egg (which I value highest), that leaves 3 apples and 1 soap. The other player might want the soap and some apples. If they take the soap and all 3 apples, that's possible. But if they also want the egg, there's a conflict. However, since I can't know their preferences, I need to structure my proposal in a way that allows the other player to take what's left without overlapping.} \\
\\ 
\texttt{Alternatively, maybe I should take the soap and some apples. If I take the soap (3 points) and, say, 1 apple (1 point), that's 4 points total. Then the other player could take the remaining 2 apples and the egg. But if the other player values the egg highly, they might want it, which would conflict if I take the soap. Wait, but if I take the soap and 1 apple, and they take the egg and 2 apples, that's possible. But my total would be 3+1=4, which is better than just taking the egg (4) alone. Hmm, same points. But maybe taking the egg is better because it's a single item, leaving more flexibility for the other player.} \\
\\ 
\texttt{Another angle: if I take the egg, the other player can take the soap and all 3 apples. That would give me 4 points, and the other player would get their value for soap and apples. If their value for soap is high, that's good for them. If they value apples more, maybe they prefer more apples. Either way, as long as their proposal doesn't require the egg, it's okay.} \\
\\ 
\texttt{But what if the other player also wants the egg? Then our proposals would conflict. But since I don't know their values, I have to assume they might not. Maybe the best strategy is to take the highest value item for} \\
            }
        }
    }
     \\ \\

    \theutterance \stepcounter{utterance}  
    & & & \multicolumn{2}{p{0.3\linewidth}}{
        \cellcolor[rgb]{0.9,0.9,0.9}{
            \makecell[{{p{\linewidth}}}]{
                \texttt{\tiny{[GM$|$GM]}}
                \texttt{continue} \\
            }
        }
    }
    & & \\ \\

    \theutterance \stepcounter{utterance}  
    & & & \multicolumn{2}{p{0.3\linewidth}}{
        \cellcolor[rgb]{0.9,0.9,0.9}{
            \makecell[{{p{\linewidth}}}]{
                \texttt{\tiny{[GM$|$GM]}}
                \texttt{<think>} \\
\texttt{Okay, let's try to figure out the best way to approach this. My values are apple:1, soap:3, egg:4. The other player's values are unknown. The key is to make a proposal that maximizes my points while ensuring that the other player's proposal doesn't conflict with mine.} \\
\\ 
\texttt{First, I need to consider what the other player might want. Since I don't know their values, I have to make assumptions. The items available are 3 apples, 1 soap, 1 egg. The other player could value any of these. If I take the egg (which I value highest), that leaves 3 apples and 1 soap. The other player might want the soap and some apples. If they take the soap and all 3 apples, that's possible. But if they also want the egg, there's a conflict. However, since I can't know their preferences, I need to structure my proposal in a way that allows the other player to take what's left without overlapping.} \\
\\ 
\texttt{Alternatively, maybe I should take the soap and some apples. If I take the soap (3 points) and, say, 1 apple (1 point), that's 4 points total. Then the other player could take the remaining 2 apples and the egg. But if the other player values the egg highly, they might want it, which would conflict if I take the soap. Wait, but if I take the soap and 1 apple, and they take the egg and 2 apples, that's possible. But my total would be 3+1=4, which is better than just taking the egg (4) alone. Hmm, same points. But maybe taking the egg is better because it's a single item, leaving more flexibility for the other player.} \\
\\ 
\texttt{Another angle: if I take the egg, the other player can take the soap and all 3 apples. That would give me 4 points, and the other player would get their value for soap and apples. If their value for soap is high, that's good for them. If they value apples more, maybe they prefer more apples. Either way, as long as their proposal doesn't require the egg, it's okay.} \\
\\ 
\texttt{But what if the other player also wants the egg? Then our proposals would conflict. But since I don't know their values, I have to assume they might not. Maybe the best strategy is to take the highest value item for} \\
            }
        }
    }
    & & \\ \\

    \theutterance \stepcounter{utterance}  
    & & \multicolumn{4}{p{0.6\linewidth}}{
        \cellcolor[rgb]{0.9,0.9,0.9}{
            \makecell[{{p{\linewidth}}}]{
                \texttt{\tiny{[P1$\langle$GM]}}
                \texttt{<think>} \\
\texttt{Okay, let's try to figure out the best way to approach this. My values are apple:1, soap:3, egg:4. The other player's values are unknown. The key is to make a proposal that maximizes my points while ensuring that the other player's proposal doesn't conflict with mine.} \\
\\ 
\texttt{First, I need to consider what the other player might want. Since I don't know their values, I have to make assumptions. The items available are 3 apples, 1 soap, 1 egg. The other player could value any of these. If I take the egg (which I value highest), that leaves 3 apples and 1 soap. The other player might want the soap and some apples. If they take the soap and all 3 apples, that's possible. But if they also want the egg, there's a conflict. However, since I can't know their preferences, I need to structure my proposal in a way that allows the other player to take what's left without overlapping.} \\
\\ 
\texttt{Alternatively, maybe I should take the soap and some apples. If I take the soap (3 points) and, say, 1 apple (1 point), that's 4 points total. Then the other player could take the remaining 2 apples and the egg. But if the other player values the egg highly, they might want it, which would conflict if I take the soap. Wait, but if I take the soap and 1 apple, and they take the egg and 2 apples, that's possible. But my total would be 3+1=4, which is better than just taking the egg (4) alone. Hmm, same points. But maybe taking the egg is better because it's a single item, leaving more flexibility for the other player.} \\
\\ 
\texttt{Another angle: if I take the egg, the other player can take the soap and all 3 apples. That would give me 4 points, and the other player would get their value for soap and apples. If their value for soap is high, that's good for them. If they value apples more, maybe they prefer more apples. Either way, as long as their proposal doesn't require the egg, it's okay.} \\
\\ 
\texttt{But what if the other player also wants the egg? Then our proposals would conflict. But since I don't know their values, I have to assume they might not. Maybe the best strategy is to take the highest value item for} \\
\\ 
\\ 
\texttt{The time is up. It is now your turn to submit a secret proposal.} \\
            }
        }
    }
    & \\ \\

    \theutterance \stepcounter{utterance}  
    & \multicolumn{4}{p{0.6\linewidth}}{
        \cellcolor[rgb]{0.9,0.9,0.9}{
            \makecell[{{p{\linewidth}}}]{
                \texttt{\tiny{[P1$\rangle$GM]}}
                \texttt{<think>} \\
\texttt{Okay, let's try to figure out the best way to approach this. My values are apple:1, soap:3, egg:4. The other player's values are unknown. The key is to make a proposal that maximizes my points while ensuring that the other player's proposal doesn't conflict with mine.} \\
\\ 
\texttt{First, I need to consider what the other player might want. Since I don't know their values, I have to make assumptions. The items available are 3 apples, 1 soap, 1 egg. The other player could value any of these. If I take the egg (which I value highest), that leaves 3 apples and 1 soap. The other player might want the soap and some apples. If they take the soap and all 3 apples, that's possible. But if they also want the egg, there's a conflict. However, since I can't know their preferences, I need to structure my proposal in a way that allows the other player to take what's left without overlapping.} \\
\\ 
\texttt{Alternatively, maybe I should take the soap and some apples. If I take the soap (3 points) and, say, 1 apple (1 point), that's 4 points total. Then the other player could take the remaining 2 apples and the egg. But if the other player values the egg highly, they might want it, which would conflict if I take the soap. Wait, but if I take the soap and 1 apple, and they take the egg and 2 apples, that's possible. But my total would be 3+1=4, which is better than just taking the egg (4) alone. Hmm, same points. But maybe taking the egg is better because it's a single item, leaving more flexibility for the other player.} \\
\\ 
\texttt{Another angle: if I take the egg, the other player can take the soap and all 3 apples. That would give me 4 points, and the other player would get their value for soap and apples. If their value for soap is high, that's good for them. If they value apples more, maybe they prefer more apples. Either way, as long as their proposal doesn't require the egg, it's okay.} \\
\\ 
\texttt{But what if the other player also wants the egg? Then our proposals would conflict. But since I don't know their values, I have to assume they might not. Maybe the best strategy is to take the highest value item for} \\
            }
        }
    }
    & & \\ \\

    \theutterance \stepcounter{utterance}  
    & & & \multicolumn{2}{p{0.3\linewidth}}{
        \cellcolor[rgb]{0.9,0.9,0.9}{
            \makecell[{{p{\linewidth}}}]{
                \texttt{\tiny{[GM$|$GM]}}
                \texttt{continue} \\
            }
        }
    }
    & & \\ \\

    \theutterance \stepcounter{utterance}  
    & & & \multicolumn{2}{p{0.3\linewidth}}{
        \cellcolor[rgb]{0.9,0.9,0.9}{
            \makecell[{{p{\linewidth}}}]{
                \texttt{\tiny{[GM$|$GM]}}
                \texttt{player was instructed to make a proposal} \\
            }
        }
    }
    & & \\ \\

\end{supertabular}
}

\end{document}
